\documentclass[10pt,a4paper]{amsart}
\usepackage[margin=19mm]{geometry}
\usepackage{amsmath,amssymb,mathtools}
\usepackage[scr=rsfs,cal=euler]{mathalfa}
\usepackage{xfrac}
\usepackage[unicode,hidelinks,bookmarks=false]{hyperref}
\newcommand{\ie}{i.e.\ }
\newcommand{\cf}{cf.\ }
\newcommand{\resp}{respectively\ }
\newcommand{\Subsection}{\subsection}
\providecommand{\nobreakdash}{\nobreak\hbox{-}}
\setlength{\emergencystretch}{2em}
\multlinegap0pt
\usepackage{bm}
\usepackage{mathtools}
\usepackage{tikz}
\usepackage[shortlabels]{enumitem}
\newtheorem{definition}{Definition}
\newtheorem{lemma}{Lemma}
\usepackage{cleveref}
\crefname{definition}{Definition}{Definitions}
\crefname{lemma}{Lemma}{Lemmas}
\newcommand{\ind}{\mathbbm{1}}
\newcommand{\mathbbm}[1]{\text{\usefont{U}{DSSerif}{m}{n}#1}}
\title{Entrywise Error Bounds for Spectral Ranking with Semi-Random Adversaries}
\author{Dongmin Lee, Anuran Makur, Japneet Singh}
\date{Source: arXiv:2605.23854v1 (CC BY 4.0)}
\begin{document}
\maketitle

\noindent\emph{Source and licence.} Entrywise Error Bounds for Spectral Ranking with Semi-Random Adversaries. Version arXiv:2605.23854v1 (CC BY 4.0), distributed by arXiv under the Creative Commons Attribution 4.0 International licence, http://creativecommons.org/licenses/by/4.0/, as recorded in the arXiv licence metadata for this exact version and shown on the abstract page https://arxiv.org/abs/2605.23854v1. Full licence notice: \texttt{licenses/arxiv-2605.23854-CC-BY-4.0.txt}.\par\smallskip

\noindent\textit{Editorial scope.} Counted unit: the article's lemma lem:recovering-properties-of-erdos-renyi-graphs (source0.tex lines 492--502). The article's complete original proof of it is delivered with it (source0.tex lines 504--517, one \texttt{proof} environment of 14 lines). No mathematics is written, repaired or re-derived: the statement and the proof are the article's own lines, in the article's own notation, and no retained line is substituted or re-flowed.  Retained uncounted as the source-local context the proof needs: lines 488--490.  Nothing was omitted from any retained span, so the package carries no omission record.  No retained line contains a citation, so the wrapper carries no bibliography and no bibliography entry.  No retained line contains a figure environment, an image-inclusion command or drawing code, so no image is delivered and the package declares no asset dependency.  Editorial formatting only, in addition: an AMS \texttt{amsart} preamble that restates the article's own declarations and macros used by the retained lines, namely the environments \texttt{definition}, \texttt{lemma} on separate counters, together with the standard packages those lines need.  The retained environments print in this document as Definition 1, Lemma 1 in order of appearance, whereas the article prints the same items with the numbering of its own full document.  The complete native source of the article as distributed in the arXiv e-print for this exact version is bundled unchanged as \texttt{sources/201243/source0.tex}.
\begin{definition}[Monotone Coupling Between Erd\H{o}s-R\'{e}nyi Graphs and Semi-Random Graphs]\label{def:maximal-coupling}
Given $n\in\mathbb{N}$, $p \in (0, 1]$, and $\{q_{ij} \in [p, 1] \mid \forall i,j\in\mathbb{N}: i<j\leq n\}$, for each pair $(i, j) \in [n]^2$ such that $i < j$, draw an independent uniformly distributed random variable $U_{ij} \sim \textsf{Unif}(0, 1)$. Construct two graphs, $\mathcal{G}_\text{E}=([n], \mathcal{E}_\text{E})$ and $\mathcal{G}_\text{S}=([n], \mathcal{E}_\text{S})$, such that $(i, j) \in \mathcal{E}_\text{E}$ if and only if $U_{ij} \leq p$ and $(i, j) \in \mathcal{E}_\text{S}$ if and only if $U_{ij} \leq q_{ij}$. The graphs $\mathcal{G}_\text{E}$ and $\mathcal{G}_\text{S}$ are realizations of the \emph{monotone coupling} between the random variables $G_\text{E}$, which is distributed according to the Erd\H{o}s-R\'{e}nyi model with edge sampling probability $p$, and $G_\text{S}$, which is distributed according to the semi-random model with edge sampling probabilities $\{q_{ij}\}$. Furthermore, $\mathbb{P}(G_\text{E} \subseteq G_\text{S})=1$. 
\end{definition}

\begin{lemma}[Recovering Properties of Erd\H{o}s-R\'{e}nyi graphs from Subgraphs of Semi-Random Graphs]\label{lem:recovering-properties-of-erdos-renyi-graphs}
Let $\mathcal{P}$ be some graph property that holds for an Erd\H{o}s-R\'{e}nyi graph (with edge sampling probability $p$) with high probability, i.e., 
\begin{equation}
\mathbb{P}_{G_\text{E} \sim \mathsf{ER}(n,p)}(G_\text{E} \in \mathcal{P}) \geq 1 - O(n^{-t})
\end{equation}
for some fixed $t>0$.
Let $G_\text{S}$ be a semi-randomly generated graph under edge sampling probabilities $\{q_{ij}\}$. Then, with high probability, there must exist a subgraph of $G_\text{S}$ that satisfies $\mathcal{P}$, i.e.,
\begin{equation}
\mathbb{P}_{G_\text{S} \sim \mathsf{SR}(n,\{q_{ij}\})}(\exists \mathcal{G} \subseteq G_\text{S}: \mathcal{G} \in \mathcal{P}) \geq 1 - O(n^{-t}).
\end{equation}
\end{lemma}

\begin{proof}
Given a graph $\mathcal{G}_\text{S}$, let $G_\text{E}$ be distributed according to $\mathbb{P}(G_\text{E}\mid G_\text{S}=\mathcal{G}_\text{S})$ under the coupling defined in \cref{def:maximal-coupling}. Note that $\mathbb{P}(G_\text{E} \subseteq \mathcal{G}_\text{S})=1$ by definition. Then,
\begin{equation}
\mathbb{P}_{G_\text{E}\mid G_\text{S}}(G_\text{E} \in \mathcal{P} \mid G_\text{S}=\mathcal{G}_\text{S}) \leq \ind\{\exists \mathcal{G} \subseteq \mathcal{G}_\text{S}: \mathcal{G} \in \mathcal{P}\}.
\end{equation}
Thus,
\begin{align}
&\phantom{{}={}}\mathbb{P}_{G_\text{S} \sim \mathsf{SR}(n,\{q_{ij}\})}(\exists \mathcal{G} \subseteq G_\text{S}: \mathcal{G} \in \mathcal{P})
\\ &= \mathbb{E}_{G_\text{S} \sim \mathsf{SR}(n,\{q_{ij}\})}[\ind\{\exists \mathcal{G} \subseteq G_\text{S}: \mathcal{G} \in \mathcal{P}\}] \\
&\geq \mathbb{E}_{G_\text{S} \sim \mathsf{SR}(n,\{q_{ij}\})}[\mathbb{P}_{G_\text{E}|G_\text{S}}(G_\text{E} \in \mathcal{P} \mid G_\text{S})] \\
&= \mathbb{P}_{G_\text{E} \sim \mathsf{ER}(n,p)}(G_\text{E} \in \mathcal{P}) \geq 1 - O(n^{-t}),
\end{align}
which completes the proof.
\end{proof}

\end{document}
